\documentclass[10pt,a4paper]{amsart}
\usepackage[margin=19mm]{geometry}
\usepackage{amsmath,amssymb,mathtools}
\usepackage[scr=rsfs,cal=euler]{mathalfa}
\usepackage{xfrac}
\usepackage[unicode,hidelinks,bookmarks=false]{hyperref}
\newcommand{\ie}{i.e.\ }
\newcommand{\cf}{cf.\ }
\newcommand{\resp}{respectively\ }
\newcommand{\Subsection}{\subsection}
\providecommand{\nobreakdash}{\nobreak\hbox{-}}
\setlength{\emergencystretch}{2em}
\multlinegap0pt
\usepackage{bm}
\usepackage{bbm}
\usepackage{dsfont}
\usepackage{xspace}
\usepackage{cancel}
\usepackage{mathtools}
\usepackage{amsthm}
\makeatletter
\@ifundefined{theorem}{\newtheorem{theorem}{Theorem}}{}
\@ifundefined{lemma}{\newtheorem{lemma}[theorem]{Lemma}}{}
\makeatother
\makeatletter
\newcommand{\N}{\mathbb{N}}
\newcommand{\R}{\mathbb{R}}
\newcommand{\bx}{\bm{x}}
\expandafter\ifx\csname claimproof\endcsname\relax
\newenvironment{claimproof}[1][Proof]{\begin{proof}[#1]\renewcommand*{\qedsymbol}{\(\diamondsuit\)}}{\end{proof}}
\fi
\newcommand{\la}{\lambda}
\newcommand{\sset}[1]{\left\{{#1}\right\}}
\newcommand{\tbx}{\tilde{\bx}}
\newcommand{\tx}{\tilde{x}}
\makeatother
\title[Forbidden induced subgraphs for graphs and signed graphs wit]{Forbidden induced subgraphs for graphs and signed graphs with eigenvalues bounded from below}
\author{Zilin Jiang, Alexandr Polyanskii}
\date{Source: arXiv:2111.10366v3 (CC BY 4.0)}
\begin{document}
\maketitle

\noindent\emph{Source and licence.} Forbidden induced subgraphs for graphs and signed graphs with eigenvalues bounded from below. Version arXiv:2111.10366v3 (CC BY 4.0), distributed under the Creative Commons Attribution 4.0 International licence, as recorded in the arXiv licence metadata for this exact version and shown on the abstract page https://arxiv.org/abs/2111.10366v3. Full rights notice: licenses/arxiv-2111.10366-CC-BY-4.0.txt.\par\smallskip

\noindent\textit{Editorial scope.} Counted unit: the Lemma whose source label is \texttt{lem:path-clique-extension} in the article (source0.tex lines 1316--1325, source label \texttt{lem:path-clique-extension}), delivered together with the article's own complete original proof of it (source0.tex lines 1327--1374, one \texttt{proof} environment of 48 lines). No mathematics is written, repaired or re-derived: statement and proof are the article's own text line for line. Required context retained uncounted: none. Every definition, display and label of those lines is retained unchanged. Omitted from this package and not redistributed: 1--1315, 1326--1326, 1375--3150. Nothing inside a retained excerpt is omitted or altered: every retained body is delivered exactly as the article's lines are written, so no omission receipt of any kind accompanies this package. Citations: the retained excerpts carry no citation command at all, so no bibliography entry of the article is transcribed and no reference is redistributed. Reference adaptation: none was needed and none was made; every reference inside the delivered unit resolves inside it, so no unresolved reference or citation is produced. Printed slips kept literal: no printed word, symbol, hypothesis or number of the counted statement or of its proof was altered. Layout and class: the article's own class is \texttt{article} with packages including hyperref, amsfonts, amsmath, amssymb, amsthm, fullpage, enumitem, tikz, bm, mathtools, tabularx, xparse, cleveref; the wrapper is the lane's pinned \texttt{amsart} editorial wrapper at 10pt on A4, which restates the theorem-like environments the retained text needs and adds the article's own macro definitions that the retained text needs (\texttt{\textbackslash N}, \texttt{\textbackslash R}, \texttt{\textbackslash bx}, \texttt{\textbackslash la}, \texttt{\textbackslash sset}, \texttt{\textbackslash tbx}, \texttt{\textbackslash tx}), with extra wrapper packages (bm, bbm, dsfont, xspace, cancel, mathtools). The delivered numbering is the unit's own. Assets: no retained span contains a figure environment, a graphics-inclusion command, a PDF-page inclusion, a table or a TikZ picture; no image file is copied into this package, the package declares no asset dependency and no image file is redistributed with it. Licence: version arXiv:2111.10366v3 of this article is distributed under the Creative Commons Attribution 4.0 International licence, as recorded in the arXiv licence metadata for this exact version and shown on the abstract page https://arxiv.org/abs/2111.10366v3. A full rights notice is delivered at licenses/arxiv-2111.10366-CC-BY-4.0.txt. Characters: every retained excerpt is pure ASCII, so the delivered engine, XeLaTeX with its default Latin Modern text font, reports no missing character.
\begin{lemma} \label{lem:path-clique-extension}
    Suppose that $A$ is a nonempty vertex subset of a graph $F$ and $\la \ge 2$. If the path extensions of $F$ satisfy
    \[
        \lim_{\ell\to\infty} \la_1(F, A, \ell) < -\la,
    \]
    then there exists $m \in \N^+$ such that the path-clique extensions of $F$ satisfy
    \[
        \la_1(F, A, \ell, K_m) < -\la \text{ for every }\ell \in \N.
    \]
\end{lemma}

\begin{proof}
    Pick $\ell_0 \in \N$ such that $\la_1(F, A, \ell) < -\la$ for every $\ell \ge \ell_0$. Clearly $\la_1(F, A, \ell, K_m) < -\la$ for every $\ell \ge \ell_0$ and every $m \in \N^+$. It suffices to show that for every $\ell \in \sset{0, \dots, \ell_0 - 1}$ there exists $m \in \N^+$ such that $\la_1(F, A, \ell, K_m) < -\la$.
    
    Let $v_0 \dots v_{\ell_0}$ be the path added to $F$ to obtain $(F, A, \ell_0)$, where the vertex $v_0$ is connected to every vertex in $A$. Set $\la_0 := -\la_1(F, A, \ell_0)$, and let $\bx \colon V(F) \cup \sset{v_0, \dots, v_{\ell_0}} \to \R$ be an eigenvector of $(F, A, \ell_0)$ associated with $-\la_0$.

    Now fix $\ell \in \N$ with $\ell < \ell_0$, and let $m \in \N^+$ be determined later. Set $V_\ell := V(F) \cup \sset{v_0, \dots, v_\ell}$, and identify the vertex set of $(F, A, \ell, K_m)$ with $V_\ell \cup V(K_m)$. We abuse notation and write $x_i$ in place of $x_{v_i}$ for $i \in \sset{\ell, \dots, \ell_0}$. Define $\tbx\colon V_\ell \cup V(K_m) \to \R$ by
    \[
        \tx_v = \begin{cases}
            x_v & \text{if }v \in V_\ell; \\
            x_{\ell+1}/m & \text{if }v \in V(K_m).
        \end{cases}
    \]
    
    We claim that $\sum_{v \in V_\ell}x_v^2 > 0$. Indeed, assume for the sake of contradiction that $x_v = 0$ for $v \in V_\ell$. Using $-\la_0x_i = \sum_{u\sim v_i}x_u$ for $i \in \sset{\ell, \dots, \ell_0-1}$, where the sum is taken over all vertices $u$ that are adjacent to the vertex $v_i$ in $(F, A, \ell_0)$, one can prove inductively that $x_i = 0$ for $i \in \sset{\ell + 1, \dots, \ell_0}$, which contradicts the assumption that $\bx$ is a nonzero vector.
    
    Because $\sum_{v \in V_\ell} x_v^2 > 0$, clearly $\tbx$ is a nonzero vector. We compute
    \[
        \tbx^\intercal \tbx = \sum_{v \in V_\ell}x_v^2 + m(x_{\ell + 1}/m)^2.
    \]
    Moreover we can compute $\tbx^\intercal A_{(F, A, \ell, K_m)} \tbx$ as follows
    \[
        \tbx^\intercal A_{(F, A, \ell, K_m)} \tbx = \sum_{u, v \in V_\ell\colon u \sim v} x_ux_v + 2x_\ell x_{\ell+1} + m(m-1)(x_{\ell+1}/m)^2.
    \]
    Since $\bx$ is an eigenvector of $(F, A, \ell_0)$ associated with $-\la_0$, we obtain that
    \[
        \sum_{u, v \in V_\ell\colon u \sim v} x_ux_v + x_\ell x_{\ell + 1} = \sum_{v \in V_\ell}x_v\sum_{u \sim v}x_u = -\la_0 \sum_{v \in V_\ell}x_v^2.
    \]
    Thus $\tbx^\intercal A_{(F, A, \ell, K_m)} \tbx$ can be simplified to
    \[
        \tbx^\intercal A_{(F, A, \ell, K_m)} \tbx = -\la_0 \sum_{v \in V_\ell}x_v^2 + x_\ell x_{\ell+1} + m(m-1)(x_{\ell+1}/m)^2.
    \]
    The Rayleigh principle says that $\la_1(F, A, \ell, K_m)$ is at most
    \[
        \frac{-\la_0 \sum_{v \in V_\ell} x_v^2 + x_\ell x_{\ell+1} + m(m-1)(x_{\ell+1}/m)^2}{\sum_{v \in V_\ell}x_v^2 + m(x_{\ell + 1}/m)^2},
    \]
    which, as $m \to \infty$, approaches
    \[
        -\la_0 + \frac{(x_\ell + x_{\ell+1})x_{\ell+1}}{\sum_{v \in V_\ell}x_v^2}.
    \]
    Here we used the above claim that the denominator in the limit is positive.
    
    Recall that $\la_0 = -\la_1(F, A, \ell_0) > \la \ge 2$. It suffices to show that $(x_\ell + x_{\ell+1})x_{\ell+1} \le 0$. In fact, we prove inductively that $(x_i + x_{i+1})x_{i+1} \le 0$ for $i \in \sset{\ell, \dots, \ell_0 - 1}$. The base case where $i = \ell_0 - 1$ follows immediately from $-\la_0x_{\ell_0} = x_{\ell_0-1}$ and $\la_0 > 2$. For the inductive step, using $-\la_0x_{i+1} = x_i + x_{i+2}$ and $\la_0 > 2$, we obtain
    \begin{multline*}
        (x_i + x_{i+1})x_{i+1} = (-\la_0x_{i+1} - x_{i+2} + x_{i+1})x_{i+1} = -(\la_0 - 2)x_{i+1}^2 - (x_{i+1}+x_{i+2})x_{i+1} \\
        \le -(x_{i+1}+x_{i+2})x_{i+1} = -(x_{i+1}+x_{i+2})^2 + (x_{i+1}+x_{i+2})x_{i+2} \le (x_{i+1}+x_{i+2})x_{i+2},
    \end{multline*}
    which is nonpositive by the inductive hypothesis.
\end{proof}

\end{document}
